\subsection{Deterministic Finite Automata (DFA)}

Secara intuitif, \textit{Deterministic Finite Automata} (DFA) dapat dipahami sebagai "mesin pembaca" yang bergerak langkah demi langkah saat menerima karakter masukan. Pada setiap langkah, mesin selalu berada tepat pada satu keadaan (\textit{state}) tertentu, lalu berpindah ke keadaan berikutnya berdasarkan karakter yang sedang dibaca. Sifat paling penting dari DFA adalah deterministik: untuk setiap pasangan keadaan dan simbol masukan, hanya ada satu perpindahan yang mungkin. Karena itu, perilaku DFA dapat diprediksi dan sangat cocok digunakan pada tahap lexical analysis yang menuntut keputusan cepat serta konsisten untuk setiap karakter yang dibaca dari kiri ke kanan \cite{hopcroft2006automata,sipser2012introduction}.

Secara formal, sebuah DFA didefinisikan sebagai \textit{quintuple} di bawah ini
\[A = (Q, \Sigma, \delta, q_0, F).\] Himpunan berhingga $Q$ menyatakan semua state yang mungkin, $\Sigma$ adalah himpunan simbol masukan, $\delta : Q \times \Sigma \rightarrow Q$ adalah fungsi transisi yang memetakan pasangan state beserta sebuah simbol masukan ke sebuah state lain, $q_0 \in Q$ adalah state awal dari mesin, dan $F \subseteq Q$ adalah himpunan \textit{acceptance state}. Untuk sebuah string masukan $w$, DFA memproses simbol demi simbol melalui perluasan fungsi transisi $\hat{\delta}$, yaitu string dinyatakan diterima jika state akhir pemrosesan berada di dalam $F$. Formulasi ini memberi dasar matematis yang kuat untuk menyatakan apakah suatu pola leksikal termasuk bahasa yang diinginkan atau tidak \cite{hopcroft2006automata,sipser2012introduction}.

Lewat quintuple ini, dapat dimodelkan sebuah mesin yang dapat menerima sejumlah input karakter $\Sigma$, dengan state awal dari mesin adalah $q_0$, yang memiliki sejumlah state mesin disimpan dalam $Q$ dan dapat berpindah-pindah dari satu state ke state yang lain menggunakan fungsi transisi $\delta$. Sejumlah input karakter merupakan sebuah "bahasa" yang valid jika state akhir dari mesin tersebut berada pada himpunan state final $F$.

Komponen pembangun DFA pada implementasi lexer dapat dipetakan langsung dari definisi formal tersebut. State merepresentasikan konteks pembacaan token, misalnya state awal, state saat membaca identifier, state saat membaca bilangan, atau state saat berada di dalam komentar. Alfabet masukan berupa karakter-karakter source code yang relevan (huruf, digit, operator, delimiter, dan simbol khusus lain). Fungsi transisi diwujudkan sebagai aturan \textit{if/switch} yang menentukan state berikutnya untuk setiap karakter yang masuk. State awal dipakai ketika belum ada token yang sedang dibangun, sedangkan state penerima menandai bahwa pola token valid telah lengkap dan token siap di-\textit{emit}. Dengan pemetaan ini, desain DFA tidak hanya teoretis, tetapi langsung menjadi kerangka arsitektur program lexer yang modular dan mudah diuji \cite{aho2007compilers,grune2012modern}.

Aplikasi DFA sangat luas pada perangkat lunak sistem, dan pada konteks tugas besar ini aplikasi utamanya adalah lexical analyzer/compiler front-end. DFA dipakai untuk mengenali lexeme seperti keyword, identifier, integer, real, operator, dan komentar secara efisien dalam waktu linear terhadap panjang input. Di luar compiler, konsep yang sama juga digunakan pada mesin pencocokan pola berbasis regular expression, penyaring teks, validasi format masukan, hingga sebagian mekanisme inspeksi paket pada jaringan yang membutuhkan klasifikasi string secara cepat dan deterministik. Karena kesederhanaan model serta jaminan performanya, DFA menjadi fondasi praktis yang menjembatani teori bahasa formal dengan implementasi perangkat lunak nyata \cite{aho2007compilers,hopcroft2006automata}.

Sebagai ilustrasi konkret, perhatikan DFA sederhana yang dirancang untuk menerima semua string yang dimulai dengan ``01'' (atas alfabet biner $\{0, 1\}$). DFA ini memiliki empat state: state awal $q_0$, state setelah membaca `0' yaitu $q_1$, state penerima $q_2$ yang dicapai setelah membaca ``01'', dan \textit{dead state} $q_3$ yang merepresentasikan penolakan (karena pola tidak cocok).

Diagram transisi DFA tersebut dapat divisualisasikan sebagai berikut.

\begin{figure}[H]
    \centering
    \begin{tikzpicture}[scale=1.2, auto, node distance=2.8cm]
        % Node definitions
        \node[state, initial] (q0) {$q_0$};
        \node[state, right of=q0] (q1) {$q_1$};
        \node[state, accepting, right of=q1] (q2) {$q_2$};
        \node[state, below of=q0] (q3) {$q_3$};

        % Transitions
        \draw (q0) edge[above] node{0} (q1);
        \draw (q0) edge[below left] node{1} (q3);
        \draw (q1) edge[above] node{1} (q2);
        \draw (q1) edge[below] node{0} (q3);
        \draw (q2) edge[loop above] node{0, 1} (q2);
        \draw (q3) edge[loop below] node{0, 1} (q3);
    \end{tikzpicture}
    \caption{Diagram transisi DFA untuk menerima string dimulai dengan ``01'' (dengan dead state)}
    \label{fig:dfa-starts-with-01}
\end{figure}

Diagram transisi di atas juga dapat direpresentasikan dengan menggunakan sebuah tabel transisi. Semua kombinasi state dan simbol masukan memiliki transisi yang jelas (sesuai dengan konsep DFA itu sendiri). Berikut adalah tabelnya.

\begin{center}
    \small
    \setlength{\extrarowheight}{2pt}
    \begin{tabular}{|c|c|c|}
        \hline
        \textbf{State}    & \textbf{0} & \textbf{1} \\
        \hline
        $\rightarrow q_0$ & $q_1$      & $q_3$      \\
        \hline
        $q_1$             & $q_3$      & $q_2$      \\
        \hline
        $q_2\star$        & $q_2$      & $q_2$      \\
        \hline
        $q_3$             & $q_3$      & $q_3$      \\
        \hline
    \end{tabular}
    \captionof{table}{Tabel transisi DFA untuk string dimulai dengan ``01''}
\end{center}

\textit{Dead state} (atau \textit{reject state}) adalah state yang bukan penerima dan dari mana mesin tidak dapat keluar menuju state penerima. Setelah mesin memasuki dead state, ia akan tetap berada di state tersebut untuk semua input yang tersisa, sehingga string otomatis ditolak. Dalam kasus ini, $q_3$ merupakan sebuah dead state. Sekali mesin tersebut masuk ke dalam state $q_3$, maka apapun masukan selanjutnya akan tetap membuat mesin berada pada $q_3$. Penggunaan dead state penting untuk menjamin bahwa fungsi transisi $\delta$ selalu terdefinisi untuk setiap kombinasi state dan simbol masukan, sesuai dengan definisi formal DFA.

Proses penerimaan string dapat diikuti langkah demi langkah. Sebagai contoh, Untuk string masukan ``\texttt{0110}'': dimulai dari $q_0$, karakter `0' membawa ke $q_1$; karakter `1' membawa ke $q_2$ (state penerima); karakter `1' tetap di $q_2$; karakter `0' tetap di $q_2$. Pada akhir string, mesin berada di state penerima $q_2$, sehingga ``\texttt{0110}'' diterima.

Contoh yang lain, Untuk string ``\texttt{10}'' yang dimulai dengan `1', mesin membaca `1' dari $q_0$ dan berpindah ke dead state $q_3$; kemudian membaca `0' dari $q_3$ dan tetap di $q_3$. Pada akhir string, mesin berada di dead state $q_3$, sehingga string ditolak.
\subsection{Token}

Token adalah unit terkecil dari source code yang bermakna dan dapat dikenali oleh lexer. Dalam konteks analisis leksikal, \textit{token} berfungsi sebagai jembatan antara rangkaian karakter mentah (source code) dan struktur sintaksis yang akan diproses oleh parser pada tahap berikutnya. Setiap token terdiri dari dua komponen utama: tipe token (\textit{token type}) dan nilai token (\textit{token value}). Tipe token mengidentifikasi kategori semantik dari token tersebut—misalnya apakah token itu keyword, identifier, operator, atau literal konstanta. Nilai token merupakan string spesifik yang dikembalikan dari source code, digunakan khususnya untuk token yang memerlukan informasi detail seperti identifier berupa ``x'' atau konstanta berupa ``42''.

Contoh sederhana dari rangkaian token adalah sebagai berikut. Jika source code berisi pernyataan ``\texttt{x := 10}'', lexer akan menghasilkan tiga token: (1) token dengan tipe IDENT dan nilai `` x'', (2) token dengan tipe BECOMES dan tanpa nilai khusus, dan (3) token dengan tipe INTCON dan nilai ``10''. Setiap token menyimpan informasi tambahan seperti nomor baris dan kolom tempat token ditemukan, membantu dalam pelaporan error yang akurat jika diperlukan pada tahap parsing atau analisis semantik.

Klasifikasi token dalam bahasa pemrograman umumnya mencakup beberapa kategori besar. Kategori pertama adalah \textit{keyword} atau kata kunci, yang merupakan simbol-simbol tereservasi dengan makna khusus dalam bahasa, seperti ``\texttt{if}'', ``\texttt{while}'', atau ``\texttt{function}''. Kategori kedua adalah \textit{identifier}, yang merupakan nama-nama yang didefinisikan oleh programmer untuk variabel, fungsi, atau tipe data, dan umumnya bersifat case-insensitive di bahasa tertentu. Kategori ketiga adalah \textit{konstanta atau literal}, yang merepresentasikan nilai tetap seperti bilangan integer (``\texttt{42}''), bilangan riil (``\texttt{3.14}''), karakter (``\texttt{'a'}''), atau string (``\texttt{'hello world'}''). Kategori keempat adalah \textit{operator}, yang merupakan simbol untuk operasi aritmatika (``\texttt{+}'', ``\texttt{-}''), logika (``\texttt{and}'', ``\texttt{or}''), atau perbandingan (``\texttt{==}'', ``\texttt{<}''). Kategori kelima adalah \textit{delimiter atau punctuation}, yaitu simbol-simbol struktural seperti tanda kurung, koma, titik koma, dan titik. Terakhir, kategori keenam mencakup \textit{komentar}, yang merupakan teks yang tidak mempengaruhi semantik program dan umumnya dibuang atau disimpan terpisah oleh lexer.

Peran token sangat penting dalam pipeline kompilasi modern. Dengan mengubah rangkaian karakter menjadi token, lexer menyederhanakan tugas parser: daripada parser harus menangani jutaan karakter individual, parser bekerja dengan puluhan atau ratusan token yang sudah terklasifikasi. Hal ini meningkatkan efisiensi dan kejelasan dari analisis sintaksis. Lebih lanjut, dengan mendeteksi error leksikal—seperti simbol yang tidak dikenal—pada tahap lexical analysis, proses kompilasi secara keseluruhan menjadi lebih robust dan memudahkan programmer untuk mengidentifikasi dan memperbaiki kesalahan \cite{aho2007compilers,hopcroft2006automata}.

\subsection{Lexical Analysis}